%##########################################################################
% CHAPTER 5: RESULTS
%##########################################################################
\chapter{Results}

\section{Chapter Overview}

This chapter presents the findings of the user study. All values marked as \textit{[INSERT]} are structured placeholders that will be populated following data collection. The statistical tests and reporting structure are fully specified; only the numerical values remain to be inserted.

\section{Participant Demographics}

A total of $N = 50$ participants completed the study (25 Experimental, 25 Control). \textit{[INSERT: exclusion details if applicable]} yielding a final sample of \textit{[INSERT: N]}. Table~\ref{tab:demographics} summarizes demographics.

\begin{table}[H]
\centering
\caption{Participant demographics [PLACEHOLDER].}
\label{tab:demographics}
\small
\begin{tabular}{@{}llll@{}}
\toprule
\textbf{Variable} & \textbf{Exp.\ ($n$=25)} & \textbf{Ctrl.\ ($n$=25)} & \textbf{Total ($N$=50)} \\
\midrule
Age: Mean (SD) & \textit{[INSERT]} & \textit{[INSERT]} & \textit{[INSERT]} \\
Gender: F / M / NB / PNS & \textit{[INSERT]} & \textit{[INSERT]} & \textit{[INSERT]} \\
Gaming hours/week: Mean (SD) & \textit{[INSERT]} & \textit{[INSERT]} & \textit{[INSERT]} \\
RPG experience (1--5): Mean (SD) & \textit{[INSERT]} & \textit{[INSERT]} & \textit{[INSERT]} \\
\bottomrule
\end{tabular}
\end{table}

No significant baseline differences between conditions were observed on age ($t$(\textit{[INSERT]}) = \textit{[INSERT]}, $p$ = \textit{[INSERT]}) or gaming hours ($t$(\textit{[INSERT]}) = \textit{[INSERT]}, $p$ = \textit{[INSERT]}), supporting the assumption of pre-intervention equivalence.

\section{Quantitative Results}

\subsection{GEQ Flow Subscale (H1)}

\begin{table}[H]
\centering
\caption{GEQ Flow subscale descriptive statistics [PLACEHOLDER].}
\label{tab:geq_flow}
\small
\begin{tabular}{@{}lllll@{}}
\toprule
\textbf{Condition} & $n$ & $M$ & $SD$ & \textbf{95\% CI} \\
\midrule
Experimental & 25 & \textit{[INSERT]} & \textit{[INSERT]} & \textit{[INSERT]} \\
Control & 25 & \textit{[INSERT]} & \textit{[INSERT]} & \textit{[INSERT]} \\
\bottomrule
\end{tabular}
\end{table}

Welch's independent-samples $t$-test: $t$(\textit{[INSERT]}) = \textit{[INSERT]}, $p$ = \textit{[INSERT]}, Cohen's $d$ = \textit{[INSERT]}.

\textit{[INSERT: Narrative interpretation of H1 result.]}

Shapiro-Wilk normality: Experimental $W$ = \textit{[INSERT]}, $p$ = \textit{[INSERT]}; Control $W$ = \textit{[INSERT]}, $p$ = \textit{[INSERT]}. Cronbach's $\alpha$ for GEQ Flow: \textit{[INSERT]}.

\subsection{GEQ Immersion Subscale (H2)}

\begin{table}[H]
\centering
\caption{GEQ Immersion subscale descriptive statistics [PLACEHOLDER].}
\label{tab:geq_immersion}
\small
\begin{tabular}{@{}lllll@{}}
\toprule
\textbf{Condition} & $n$ & $M$ & $SD$ & \textbf{95\% CI} \\
\midrule
Experimental & 25 & \textit{[INSERT]} & \textit{[INSERT]} & \textit{[INSERT]} \\
Control & 25 & \textit{[INSERT]} & \textit{[INSERT]} & \textit{[INSERT]} \\
\bottomrule
\end{tabular}
\end{table}

Welch's $t$-test: $t$(\textit{[INSERT]}) = \textit{[INSERT]}, $p$ = \textit{[INSERT]}, Cohen's $d$ = \textit{[INSERT]}.

\textit{[INSERT: Narrative interpretation of H2 result. Cronbach's $\alpha$ for GEQ Immersion.]}

\subsection{Personalization Perception (H3)}

\begin{table}[H]
\centering
\caption{NPS-3 Personalization Scale descriptive statistics [PLACEHOLDER].}
\label{tab:nps3}
\small
\begin{tabular}{@{}lllll@{}}
\toprule
\textbf{Condition} & $n$ & $M$ & $SD$ & \textbf{95\% CI} \\
\midrule
Experimental & 25 & \textit{[INSERT]} & \textit{[INSERT]} & \textit{[INSERT]} \\
Control & 25 & \textit{[INSERT]} & \textit{[INSERT]} & \textit{[INSERT]} \\
\bottomrule
\end{tabular}
\end{table}

Welch's $t$-test: $t$(\textit{[INSERT]}) = \textit{[INSERT]}, $p$ = \textit{[INSERT]}, Cohen's $d$ = \textit{[INSERT]}.

\textit{[INSERT: Narrative interpretation of H3 result. NPS-3 Cronbach's $\alpha$ and item-level means.]}

\subsection{Behavioural Metrics}

\begin{table}[H]
\centering
\caption{Behavioural metrics comparison [PLACEHOLDER].}
\label{tab:behavioural}
\small
\begin{tabular}{@{}llllll@{}}
\toprule
\textbf{Metric} & \textbf{Exp.\ $M$ ($SD$)} & \textbf{Ctrl.\ $M$ ($SD$)} & $t$ & $p$ & $d$ \\
\midrule
Voluntary NPC interactions & \textit{[INSERT]} & \textit{[INSERT]} & \textit{[INSERT]} & \textit{[INSERT]} & \textit{[INSERT]} \\
Dialogue read ratio & \textit{[INSERT]} & \textit{[INSERT]} & \textit{[INSERT]} & \textit{[INSERT]} & \textit{[INSERT]} \\
Exploration coverage & \textit{[INSERT]} & \textit{[INSERT]} & \textit{[INSERT]} & \textit{[INSERT]} & \textit{[INSERT]} \\
miniPXI total score & \textit{[INSERT]} & \textit{[INSERT]} & \textit{[INSERT]} & \textit{[INSERT]} & \textit{[INSERT]} \\
\bottomrule
\end{tabular}
\end{table}

\textit{[INSERT: Exploratory analysis interpretation.]}

\section{Qualitative Results (Thematic Analysis)}

Semi-structured interviews were conducted with \textit{[INSERT: $n$]} participants. Following the six-phase Reflexive Thematic Analysis procedure \citep{braun2022thematic}, \textit{[INSERT: $n$]} initial codes were generated and refined into \textit{[INSERT: $n$]} final themes.

\begin{table}[H]
\centering
\caption{Qualitative themes and representative quotations [PLACEHOLDER].}
\label{tab:themes}
\small
\begin{tabular}{@{}p{3cm}p{4cm}p{6cm}@{}}
\toprule
\textbf{Theme} & \textbf{Description} & \textbf{Representative Quote} \\
\midrule
\textit{[INSERT]} & \textit{[INSERT]} & \textit{[INSERT]} \\
\textit{[INSERT]} & \textit{[INSERT]} & \textit{[INSERT]} \\
\textit{[INSERT]} & \textit{[INSERT]} & \textit{[INSERT]} \\
\textit{[INSERT]} & \textit{[INSERT]} & \textit{[INSERT]} \\
\bottomrule
\end{tabular}
\end{table}

\section{Summary of Findings}

\textit{[INSERT: A narrative summary of the overall pattern of results across quantitative and qualitative strands, addressing whether H1, H2, and H3 were supported, what the behavioural metrics indicate, and the dominant qualitative themes.]}
